\chapter{Thyroxine (T4)}

\section*{What Is This Test?}
Thyroxine (T4, or tetraiodothyronine) is the primary hormone produced by the thyroid gland. It is synthesized in thyroid follicular cells from tyrosine residues within thyroglobulin through iodination and coupling reactions. T4 is released into the bloodstream, where over 99.97\% is bound to carrier proteins: thyroxine-binding globulin (TBG, 70\%), transthyretin (TTR, 10--15\%), and albumin (15--20\%). Only the unbound (free) fraction (FT4, approximately 0.03\% of total) is biologically active. The total T4 test measures both protein-bound and free T4. The free T4 test measures only the biologically active unbound hormone and is preferred clinically because it is not affected by changes in binding protein concentrations. T4 serves as a prohormone for T3; peripheral tissues convert T4 to active T3 through 5'-deiodination by type 1 and type 2 deiodinase enzymes. The test helps evaluate thyroid function and is most useful when combined with TSH measurement.

\section*{Alternative Names}
Total T4; Total Thyroxine; Free T4 (FT4); Free Thyroxine; Tetraiodothyronine; T4 Radioimmunoassay; T4 (Thyroxine) Free

\section*{Principle}
Free T4 is measured using a competitive immunoassay on automated analyzers. In free T4 assays, an anti-T4 antibody competes for T4 between the free T4 in the patient sample and a labeled T4 analog (chemiluminescent tracer). The labeled T4 analog does not bind significantly to albumin or TBG (thyroxine-binding globulin). The amount of labeled analog bound to the antibody is inversely proportional to the free T4 concentration. Two-step or one-step competitive immunoassay formats are used. Common platforms include: Roche Cobas ECLIA (electrochemiluminescence immunoassay using ruthenium-labeled T4 and streptavidin-coated microparticles), Abbott Architect CMIA (chemiluminescent microparticle immunoassay with acridinium-labeled T4), Siemens ADVIA Centaur/Atellica CLIA, and Beckman Coulter DxI. Equilibrium dialysis followed by liquid chromatography-tandem mass spectrometry (LC-MS/MS) is the reference method for free T4 but is not used in routine clinical practice due to complexity and cost. Total T4 is measured by a similar competitive immunoassay but with a displacement agent that releases T4 from binding proteins. Calibration is standardized against certified reference material.

\section*{History}
T4 was first crystallized by Edward Kendall in 1914 at the Mayo Clinic on Christmas Day, winning him the Nobel Prize in Physiology or Medicine (shared with Reichstein and Hench for cortisone). The first serum T4 measurement was the protein-bound iodine (PBI) test developed in the 1940s, which measured iodine content of serum proteins but was nonspecific and interfered with by dietary iodine and radiocontrast agents. The first competitive protein-binding assay for T4 was developed by Murphy and Pattee in 1964. Radioimmunoassay (RIA) for T4 was introduced by Chopra in 1972. Free T4 measurement became possible in the 1980s with labeled-antibody methods and equilibrium dialysis. Modern chemiluminescent free T4 immunoassays were introduced in the 1990s, and LC-MS/MS reference methods were standardized in the 2000s. The American Thyroid Association guidelines (2012--2017) affirmed that TSH is the primary screening test and free T4 is the first confirmatory test.

\section*{Why This Test Is Ordered}
\begin{itemize}
  \item Confirmatory test when TSH is abnormal (elevated or suppressed) to distinguish subclinical from overt thyroid dysfunction
  \item Monitoring levothyroxine replacement therapy dose adequacy in hypothyroid patients (target FT4 in upper half of reference range)
  \item Evaluation of suspected central hypothyroidism (pituitary or hypothalamic disease) where TSH is inappropriately normal or low with low FT4
  \item Assessment of thyroid status in pregnant women --- FT4 should be maintained in trimester-specific reference ranges
  \item Diagnosis of T4 thyrotoxicosis (elevated FT4 with suppressed TSH) in Graves disease or toxic multinodular goiter
  \item Monitoring amiodarone-induced thyroid dysfunction (both hypo- and hyperthyroidism) --- amiodarone inhibits T4-to-T3 conversion, so FT4 may be elevated even in euthyroid patients
  \item Workup of goiter or thyroid nodules with abnormal TSH
  \item Evaluation of suspected thyroid hormone resistance where both TSH and FT4 are elevated
  \item Neonatal screening for congenital hypothyroidism (total T4 in heel-prick blood, reflex TSH)
  \item Serial monitoring during anti-thyroid drug therapy (methimazole, propylthiouracil) for hyperthyroidism
  \item Assessment of thyroid function in patients on drugs that alter TBG levels (estrogen, tamoxifen, oral contraceptives raise TBG; androgens, glucocorticoids lower TBG)
\end{itemize}

\section*{Primary vs Secondary Test}
This is a secondary test. Free T4 is used to confirm and quantify the severity of thyroid dysfunction first detected by an abnormal TSH. It is not recommended as a standalone screening test for thyroid disease. The American Thyroid Association and Endocrine Society recommend TSH as the initial screening test, with reflex free T4 testing when TSH is abnormal. Total T4 is less clinically useful than free T4 because total T4 is affected by TBG concentration changes and is considered a third-line measurement.

\section*{How This Test Is Performed}
A healthcare professional draws a blood sample from a vein in your arm into a serum separator tube (gold-top) or plain red-top tube. Approximately 3--5 mL of blood is collected. The specimen is centrifuged to separate serum. Free T4 in serum is stable at room temperature for 24 hours, at 2--8 degrees C for 7 days, and at -20 degrees C for 6 months. Avoid repeated freeze-thaw cycles. The specimen should not be collected from an arm receiving intravenous fluids.

\section*{What Preparation Is Needed}
No fasting is required for free T4 testing. However, if the test is being used to adjust levothyroxine dosage, the blood sample should be drawn before the morning dose (trough level), since free T4 transiently rises 1--2 hours after levothyroxine ingestion. Patients on high-dose biotin supplements (>5 mg/day) should discontinue biotin 48 hours before testing because biotin interferes with streptavidin-biotin immunoassay platforms, causing falsely elevated free T4 results. Patients who have received recent intravenous iodine contrast agents should wait 4--6 weeks before testing if possible, as the iodine load transiently alters thyroid function. Heparin administration causes in-vitro lipolysis and release of free fatty acids that can displace T4 from binding proteins, artifactually elevating free T4 --- blood should be drawn before heparin administration if possible.

\section*{Contraindications and Precautions}
\begin{itemize}
  \item No absolute contraindications. Factors affecting free T4 results: pregnancy (increased TBG from estrogen elevates total T4, but free T4 remains in normal range; trimester-specific reference ranges should be used); critical illness/euthyroid sick syndrome (free T4 may be low, normal, or elevated depending on assay method); heparin therapy (in-vitro artifact increases free T4); severe non-thyroidal illness; amiodarone therapy (free T4 often in upper normal or elevated range even in euthyroid patients due to inhibited T4-to-T3 conversion); anti-thyroid antibodies and heterophile antibodies may interfere with immunoassays causing spuriously elevated results; biotin interference (falsely high free T4 in competitive assays). Drugs that increase TBG (estrogen, tamoxifen, oral contraceptives, methadone, 5-fluorouracil) will increase total T4 but not free T4. Drugs that decrease TBG (androgens, anabolic steroids, glucocorticoids) will decrease total T4 but not free T4. Drugs that displace T4 from TBG (furosemide at high doses, salicylates, phenytoin) may transiently lower total T4. Hemolysis and severe lipemia should be avoided as they affect optical detection in immunoassays.
\end{itemize}
\section*{Risks}
Minimal risk. Slight pain or bruising at venipuncture site. Rarely: hematoma, vasovagal syncope, infection, or nerve injury. No radiation exposure.

\section*{What Happens After the Test}
Results are available within 1--4 hours on standard automated analyzers. The healthcare provider interprets free T4 in the context of TSH, clinical symptoms, and physical examination. Abnormal results prompt treatment initiation or adjustment: levothyroxine therapy for hypothyroidism (starting dose 1.6 mcg/kg/day, adjusted every 6--8 weeks based on TSH and FT4); anti-thyroid drugs (methimazole 5--30 mg/day or propylthiouracil 50--150 mg three times daily) for hyperthyroidism, with dose adjustment every 4--6 weeks. During pregnancy, free T4 and TSH should be monitored every 4 weeks until mid-gestation and at least once each trimester thereafter.

\section*{Understanding Results}

\subsection*{Below Normal}
\begin{itemize}
  \item \textbf{Low free T4 with elevated TSH:} Overt primary hypothyroidism --- most commonly Hashimoto thyroiditis (anti-TPO positive), post-radioactive iodine ablation, post-thyroidectomy, or post-external beam neck radiation. Clinical manifestations include fatigue, weight gain, cold intolerance, constipation, dry skin, bradycardia, delayed deep tendon reflexes, and myxedema in severe cases.
  \item \textbf{Low free T4 with inappropriately normal or low TSH:} Central hypothyroidism (secondary or tertiary) due to pituitary adenomas (non-functioning macroadenoma compressing thyrotrope cells), Sheehan syndrome (postpartum pituitary necrosis from obstetric hemorrhage), traumatic brain injury, pituitary surgery/radiation, infiltrative diseases (sarcoidosis, hemochromatosis), or lymphocytic hypophysitis.
  \item \textbf{Low free T4 in euthyroid sick syndrome:} Critically ill patients with decreased T4-to-T3 conversion and suppressed TSH. Free T4 may be artifactually low due to free fatty acid interference. This does not typically require thyroid hormone replacement.
  \item \textbf{Antiepileptic drugs (phenytoin, carbamazepine, phenobarbital):} Increase hepatic clearance of T4, resulting in low free T4. TSH is usually normal, distinguishing from true hypothyroidism.
  \item \textbf{Severe iodine deficiency:} Endemic goiter and hypothyroidism from inadequate dietary iodine intake, particularly in mountainous regions worldwide.
\end{itemize}

\subsection*{Above Normal}
\begin{itemize}
  \item \textbf{Elevated free T4 with suppressed TSH (<0.1 mIU/L):} Overt primary hyperthyroidism. Most commonly Graves disease (autoimmune TSH receptor stimulation, diffuse goiter, orbitopathy in 25--30\%, pretibial myxedema). Free T4 is elevated above the reference range, free T3 may be disproportionately elevated.
  \item \textbf{Toxic multinodular goiter (Plummer disease):} Multiple autonomously hyperfunctioning thyroid nodules producing excess T4 and T3, most common in older adults with long-standing nontoxic goiter. TSH suppressed, elevated free T4.
  \item \textbf{Toxic adenoma (solitary hyperfunctioning nodule):} Single autonomous nodule producing excess T4/T3 with suppressed TSH. Typically occurs in younger patients (<40 years).
  \item \textbf{Subacute granulomatous thyroiditis (de Quervain thyroiditis):} Viral/post-viral inflammation causing follicular destruction and release of preformed T4/T3. Painful thyroid, elevated ESR/CRP, low radioactive iodine uptake. Elevated free T4 with suppressed TSH in thyrotoxic phase, then may progress to hypothyroid phase.
  \item \textbf{Factitious thyrotoxicosis:} Surreptitious levothyroxine ingestion (patient may deny). Elevated free T4, suppressed TSH, low thyroglobulin (endogenous thyroid suppressed), low or absent radioactive iodine uptake.
  \item \textbf{Amiodarone-induced thyrotoxicosis Type 1:} Excess iodine load drives increased thyroid hormone synthesis in abnormal thyroid (nodular goiter, latent Graves disease). Type 2: Destructive thyroiditis from direct amiodarone toxicity to follicular cells. Both present with elevated free T4, suppressed TSH. Distinguishing Type 1 from Type 2 is critical for management (anti-thyroid drugs vs. glucocorticoids).
  \item \textbf{Elevated free T4 with inappropriately normal or elevated TSH:} TSH-secreting pituitary adenoma (thyrotropinoma) or thyroid hormone resistance (THR-beta mutation). Requires pituitary MRI, alpha-subunit measurement, and genetic testing to differentiate.
  \item \textbf{Struma ovarii:} Ectopic thyroid tissue in ovarian teratoma producing excess T4/T3 with suppressed TSH. Diagnosed by whole-body radioactive iodine uptake scan showing pelvic uptake.
\end{itemize}

\section*{Normal Results}
\begin{itemize}
  \item \textbf{Free T4 (Adults):} 0.8--1.8 ng/dL (10--23 pmol/L) --- varies by laboratory and assay methodology
  \item \textbf{Free T4 (First trimester pregnancy):} 0.9--1.7 ng/dL (or trimester-specific range)
  \item \textbf{Free T4 (Second trimester pregnancy):} 0.7--1.4 ng/dL
  \item \textbf{Free T4 (Third trimester pregnancy):} 0.6--1.3 ng/dL
  \item \textbf{Free T4 (Neonates, 1--4 days):} 2.0--5.0 ng/dL (26--64 pmol/L) --- transient neonatal T4 surge
  \item \textbf{Free T4 (Children 1--5 years):} 0.8--2.0 ng/dL
  \item \textbf{Free T4 (Children 6--12 years):} 0.8--1.8 ng/dL
  \item \textbf{Total T4 (Adults):} 5.0--12.0 mcg/dL (64--155 nmol/L)
  \item \textbf{Total T4 (Pregnancy):} 7.5--18.0 mcg/dL (due to estrogen-induced TBG elevation)
  \item \textbf{Target FT4 on levothyroxine therapy:} Upper half of reference range (typically 1.2--1.8 ng/dL)
\end{itemize}

\section*{Follow-up Tests}
\begin{itemize}
  \item \textbf{TSH (third-generation):} Always interpreted alongside free T4. The TSH-free T4 relationship (log-linear) is fundamental to thyroid function interpretation. TSH is the more sensitive indicator of thyroid dysfunction.
  \item \textbf{Free T3 (FT3):} Added if hyperthyroidism is suspected with normal FT4 and suppressed TSH, or to distinguish T3 toxicosis. Also useful in monitoring anti-thyroid drug therapy.
  \item \textbf{Reverse T3 (rT3):} Elevated in non-thyroidal illness syndrome (euthyroid sick) when deiodination pathway shifts from T3 production to reverse T3 production. High rT3 with low or normal T3 and normal TSH distinguishes sick euthyroid from true central hypothyroidism.
  \item \textbf{Anti-TPO and anti-thyroglobulin antibodies:} Confirming autoimmune etiology (Hashimoto thyroiditis or Graves disease) in patients with abnormal TSH and FT4.
  \item \textbf{TSH receptor antibodies (TRAb/TSI):} Diagnostic for Graves disease when cause of hyperthyroidism is unclear. Essential in pregnancy to predict neonatal hyperthyroidism risk.
  \item \textbf{Thyroid ultrasound with Doppler:} Evaluating thyroid structure, nodularity, and vascularity. Increased vascular flow on Doppler suggests Graves disease.
  \item \textbf{Radioactive iodine uptake (I-123) and scan:} Distinguishes high-uptake states (Graves, toxic nodular goiter) from low-uptake states (thyroiditis, exogenous hormone). Performed in hyperthyroid patients when diagnosis is unclear.
  \item \textbf{Thyroglobulin:} Useful in distinguishing exogenous thyrotoxicosis (low thyroglobulin) from endogenous hyperthyroidism (elevated thyroglobulin).
  \item \textbf{Pituitary MRI with gadolinium contrast:} If central hypothyroidism or TSH-secreting adenoma is suspected based on discordant TSH and free T4 results (low TSH + low FT4, or elevated TSH + elevated FT4).
  \item \textbf{Alpha-subunit of glycoprotein hormones:} Elevated molar ratio of alpha-subunit to TSH in TSH-secreting adenomas.
\end{itemize}

\section*{References}
\begin{enumerate}
  \item Garber JR, Cobin RH, Gharib H, et al. Clinical practice guidelines for hypothyroidism in adults. Thyroid. 2012;22(12):1200-1235.
  \item Ross DS, Burch HB, Cooper DS, et al. 2016 American Thyroid Association guidelines for diagnosis and management of hyperthyroidism. Thyroid. 2016;26(10):1343-1421.
  \item Jonklaas J, Bianco AC, Bauer AJ, et al. Guidelines for the treatment of hypothyroidism. Thyroid. 2014;24(12):1670-1751.
  \item Alexander EK, Pearce EN, Brent GA, et al. 2017 Guidelines of the American Thyroid Association for the diagnosis and management of thyroid disease during pregnancy and the postpartum. Thyroid. 2017;27(3):315-389.
  \item Baloch Z, Carayon P, Conte-Devolx B, et al. Laboratory medicine practice guidelines: laboratory support for the diagnosis and monitoring of thyroid disease. Thyroid. 2003;13(1):3-126.
  \item Thienpont LM, Van Uytfanghe K, Beastall G, et al. Report of the IFCC Working Group for Standardization of Thyroid Function Tests. Clin Chem. 2010;56(6):902-911.
  \item Burtis CA, Ashwood ER, Bruns DE. Tietz Textbook of Clinical Chemistry and Molecular Diagnostics. 6th ed. Elsevier; 2023.
\end{enumerate}

\clearpage
\section*{Regulatory Status and Authorization Date}
This chapter describes a test category, not a specific commercial product. FDA, CDSCO, EU IVDR, and MHRA approval or registration dates therefore cannot be assigned without the assay manufacturer, product name, instrument, and intended use. For laboratory-developed tests, an individual agency approval date may not exist. See the Preface for the interpretation of regulatory status.

\clearpage
